\section{Limitations and Future Work}
\label{sec:limitations_future_work}

Each F1 configuration was trained once and assessed on a single fixed split. This design is appropriate for a controlled comparison, but it does not measure variation from initialization, minibatch order, or stochastic transformations. In particular, the 0.0025 Macro AP-rIoU separation between Base 1 and Base 2 is small enough that its ordering could change in repeated experiments. Subsequent work should train several seeds and accompany overall and per-class scores with confidence intervals or bootstrap estimates.

The current evidence also comes from one dataset, one detector family, and one implementation of the final evaluator. Results obtained from these cameras should not be assumed to hold for different cities, viewing angles, rotated detectors, or composition pipelines. Evaluation across cameras and external datasets would clarify whether the synthetic variants improve resilience to distribution shift despite their lower score on the present validation set.

Only one fixed synthetic release was examined. A more complete design should vary the number of inserted objects, class allocation, scale, overlap, placement region, scene compatibility, mask quality, blending procedure, and the ratio between real and generated samples. Because Motocicleta is consistently the hardest category, targeted sampling of realistic small-object examples is a priority, provided that scene density and contextual plausibility are preserved. Additional experiments should isolate LaMa correction, SAM-based extraction, placement constraints, and sampling ratio instead of changing several factors simultaneously.

Aggregate AP does not fully explain the observed rankings. Further analysis should include precision--recall curves, category confusions, duplicate predictions, localization failures, and stratification by object size, camera, lighting, congestion, and occlusion. Visual review of matched and unmatched oriented boxes could reveal why Base 1 limits false positives and why Base 2 performs better at rIoU 0.80.

Finally, data modification should be tested alongside changes to the detector itself. Promising comparisons include class-balanced objectives, rare-class resampling, larger input resolution, losses tailored to oriented boxes, hard-example mining, and alternative rotated-detection architectures. Deployment experiments should additionally report latency, memory, and energy use on the intended platform. Such measurements are necessary to determine whether the 98.93 FPS achieved by Improvement C justifies its lower accuracy in a real surveillance system.
